\subsection{Vehicle Configuration in the Right-Hand Panel}
\label{subsec:vehicle-configuration-panel}

The Vehicle tab of the right-hand configuration panel is used to build a reusable, vertically controlled rocket. It ranges from a small single-engine test vehicle to a Falcon~9-inspired booster with nine engines, grid fins, and a reaction control system (RCS). The simulator is therefore configurable rather than an exact digital copy of one rocket. Read-only values---dry mass, fuel capacity, starting fuel, aerodynamic areas, total thrust, fuel flow, and fin area---show the main physical consequences of each change.

Two immutable presets provide reproducible starting points. \emph{Falcon 9} uses the reference octaweb, four fins, RCS, and a centre-plus-two active-engine group. \emph{Simple} uses one engine without fins or RCS and faster, training-friendly timing. Any manual edit marks the build as \emph{Custom}. Users can save any number of named vehicle presets, load them, overwrite their own entries, and delete entries they no longer need. Built-in presets cannot be overwritten or deleted. A saved vehicle contains only hardware configuration; it never changes the task, reward function, trainer, environment, or telemetry settings.

The editable controls are summarized below. Geometry and mass determine how the rocket moves; actuator limits determine how quickly it can correct an error.

\begin{longtable}{>{\raggedright\arraybackslash}p{0.21\textwidth} >{\raggedright\arraybackslash}p{0.73\textwidth}}
\toprule
\textbf{Control group} & \textbf{Available selections and sliders} \\
\midrule
\endfirsthead
\toprule
\textbf{Control group} & \textbf{Available selections and sliders} \\
\midrule
\endhead
Body and fuel & \textbf{Radius} and \textbf{Height} scale the body, dry mass, tank capacity, and aerodynamic areas. Recommended fuel can follow the selected task; otherwise, \textbf{Starting Fuel} selects 1--100\% of capacity. \\
\addlinespace
Engine arrangement & Select \textbf{single}, \textbf{triple}, or nine-engine \textbf{octaweb}. Octaweb burn groups activate the centre, centre plus two, centre plus four, outer eight, or all nine engines; inactive engines remain installed mass. \textbf{Independent Engine Control} provides separate commands instead of one shared channel. \\
\addlinespace
Main engines & \textbf{Thrust per Engine} sets maximum force; \textbf{Minimum Throttle}, the lowest non-zero setting; and \textbf{Layout Spacing}, engine distance from the centreline. \textbf{Specific Impulse} represents fuel efficiency, and \textbf{Gimbal Range} limits thrust-vector tilt. \\
\addlinespace
Actuator dynamics & \textbf{Throttle Response} and \textbf{Gimbal Slew} limit command-following speed. \textbf{Startup Delay}, \textbf{Minimum Run Time}, \textbf{Restart Cooldown}, and \textbf{Shutdown State Delay} define ignition and shutdown timing. \\
\addlinespace
Grid fins & Fins can be disabled or arranged as three at $120^{\circ}$, four in ``+'', or four in asymmetric ``X''. \textbf{Radial Length}, \textbf{Tangential Width}, and \textbf{Thickness} set geometry; \textbf{Maximum Angle}, \textbf{Slew Rate}, and \textbf{Aerodynamic Effectiveness} set deflection, speed, and steering force. \\
\addlinespace
Reaction control system & RCS can be disabled; otherwise, it uses two four-nozzle pods. \textbf{Thrust per Jet}, \textbf{Specific Impulse}, and \textbf{Minimum Pulse} set force, efficiency, and shortest firing. \textbf{Nitrogen Mass} sets propellant, and \textbf{RCS Dry Mass} represents its hardware. \\
\bottomrule
\end{longtable}

Each engine uses four states. A positive command moves it from \emph{Off} to \emph{Starting} after any restart cooldown; no thrust is produced during the startup delay. In \emph{Running}, throttle is at least the configured minimum, and shutdown is blocked until the minimum run time ends. The engine then enters \emph{Shutdown}, targets zero throttle, and returns to \emph{Off} after the shutdown delay and spool-down. This is a timing model rather than a detailed combustion model, but it makes the controller anticipate actuator delay.

Vehicle changes alter mass, inertia, aerodynamics, fuel use, and control authority. The simulator builds the learning interface from the selected hardware in the background: each commandable engine group contributes throttle and two-axis gimbal actions plus its actuator-state observations; enabled fins and RCS jets contribute their own channels. General flight-state and landing-contact observations remain available to every vehicle. Disabled or unavailable hardware therefore creates no padded observation or no-op action slots. Network depth and width remain independent user-controlled trainer settings.
